\chapter*{Abstract}
\addcontentsline{toc}{chapter}{Abstract}

This Bachelor's thesis presents \textit{AI Auto Doctor}, a cross-platform mobile application for Android and iOS that addresses a fundamental gap in accessible vehicle diagnostics: the inability of existing solutions to provide verified, sensor-grounded fault diagnoses.
The system implements a two-stage diagnostic pipeline in which a deterministic Rule-Based engine (\texttt{DiagnosticEngine}) first evaluates a Diagnostic Trouble Code (DTC) against real-time OBD-II sensor data---including Short-Term Fuel Trim (STFT), Long-Term Fuel Trim (LTFT), Mass Airflow Rate (MAF), O2 sensor voltage, coolant temperature, and FreezeFrame snapshot---and produces probability-ranked hypotheses with a confidence score ranging from 40\% to 95\%.
Only after this structured verification does the system invoke Google Gemini 1.5 Flash to generate a plain-language explanation, repair cost estimate, and a deterministic ``Can Drive?'' safety verdict.

The thesis demonstrates through a controlled comparative experiment on three DTC scenarios (P0171, P0300, P0420) that the two-stage approach increases cause identification accuracy from 61\,\% (AI-only) to 87\,\%, result consistency from 58\,\% to 94\,\%, and user explainability rating from 3.2/5 to 4.5/5.
All ten functional requirements passed User Acceptance Testing on both Android and iOS.
The DiagnosticEngine executes in under 2\,ms, confirming that the rule-based layer adds negligible latency while substantially improving diagnostic reliability.

\medskip
\textbf{Keywords:} OBD-II, DTC, vehicle diagnostics, Rule-Based engine,
two-stage analysis, Flutter, Gemini AI, Bluetooth Low Energy,
Health Score, FreezeFrame, confidence scoring.

\clearpage

% ── Latvian abstract (anotacija) required by RNU
\chapter*{Anotācija}
\addcontentsline{toc}{chapter}{Anotācija}

Šis bakalaura darbs prezentē \textit{AI Auto Doctor}---krosplatformu mobilo lietotni Android un iOS, kas risina kritisku trūkumu pieejamas automobiļu diagnostikas jomā: esošo risinājumu nespēju nodrošināt verificētus, sensoru datos balstītus kļūdu diagnostikas secinājumus.
Sistēma ievieš divposmu diagnostikas konveijeru, kurā determinisks uz noteikumiem balstīts dzinējs (\texttt{DiagnosticEngine}) vispirms novērtē DTC kodu pret reāllaika OBD-II sensoru datiem---ieskaitot STFT, LTFT, MAF, O2 sensora spriegumu, dzesēšanas šķidruma temperatūru un FreezeFrame momentuzņēmumu---un tikai pēc tam aktivizē Google Gemini 1.5 Flash, lai ģenerētu skaidrojumu dabiskās valodas formā.

Kontrolēts salīdzinošais eksperiments uz trim DTC scenārijiem (P0171, P0300, P0420) uzrāda, ka divposmu pieeja palielina cēloņu identifikācijas precizitāti no 61\,\% līdz 87\,\%, rezultātu konsekvenci no 58\,\% līdz 94\,\%, un lietotāju novērtēto saprotamību no 3,2/5 līdz 4,5/5.
Visi desmit funkcionālie prasības ir izturētas.

\medskip
\textbf{Atslēgvārdi:} OBD-II, DTC, automobiļu diagnostika, uz noteikumiem balstīts dzinējs,
divposmu analīze, Flutter, Gemini AI, BLE.

\clearpage
